\section{Chronological spectral products and coarse interpolation}
\label{app:functional-details}
This appendix supplies explicit product and interpolation estimates for
Section~\ref{sec:functional-gaussian}. It also separates the treatment
of short increments from the long-block spectral argument.

\begin{lemma}[Products of differently twisted collision powers]
\label{lem:chronological-product}
Fix an integer $q\ge1$. Use one local collision space, a common smoothing
scale $\delta$, and real vectors $z_1,\ldots,z_q$ with
$|z_j|\le c\delta^2$. Let $m_j\ge1$ and assume
\[
 \delta^{-6}\sum_{j=1}^q m_j|z_j|^3\le c.
\]
Write $L_j=\mathcal L_{R,\delta,z_j}$ and let $a$ be a complex initial
insertion with bounded supremum and variation norms. Then
\begin{align}
 &\left|\ell_j\bigl(L_q^{m_q}\cdots L_1^{m_1}
                   (\mathsf S_\delta a)\nu\bigr)
       -\left(\int a\dd\nu\right)
       \exp\left(-\tfrac12\sum_{j=1}^q
                      m_j z_j^{\mathsf T}\Gamma_{R,\delta}z_j\right)
        \right|\notag\\
 &\quad\le C_q\left\{
       \delta^{-4}\sum_{j=1}^q|z_j|
       +\delta^{-6}\sum_{j=1}^qm_j|z_j|^3
       +\delta^{-2}\sum_{j=1}^q\rho^{m_j}\right\}.
       \label{eq:chronological-product-bound}
\end{align}
The index on $\ell_j$ names the chosen local space and is fixed in the
formula. The powers are in physical chronological order: the earliest
increment acts first, on the right.
\end{lemma}
\begin{proof}
Write $L_j^{m_j}=\lambda_j^{m_j}\Pi_j+E_j$.
Lemma~\ref{lem:smooth-collision-expansion} gives
$\|\Pi_j\|\le C$, $\|E_j\|\le C\rho^{m_j}$ and
$\|\Pi_j-\Pi_R\|\le C\delta^{-2}|z_j|$.
The real-frequency quadratic terms in $m_j\log\lambda_j$ are
nonpositive. The assumption on the sum of the cubic remainders bounds
each principal scalar power and every product of these powers by a
fixed constant. Expand the product into its $2^q$ terms. Every term
with a complementary factor is bounded, after evaluation, by a
constant times $\delta^{-2}\rho^{m_j}$ for at least one of its
complementary indices. The factor $\delta^{-2}$ comes only from the
initial vector, through \eqref{eq:smoothed-density-vector}; the other
projectors and powers have uniform bounds. Summing the finitely many
terms gives the last term of \eqref{eq:chronological-product-bound}.

For the all-principal term, telescope the product of projectors about
$\Pi_R$, retaining their order. Since $\Pi_R^2=\Pi_R=\nu\ell_j$,
\[
 \|\Pi_q\cdots\Pi_1-\Pi_R\|
       \le C_q\delta^{-2}\sum_j|z_j|.
\]
Evaluation on the initial vector gives the first error in
\eqref{eq:chronological-product-bound}. Its zero-frequency amplitude
is $\int a\dd\nu$. The logarithmic expansion and
$|e^w-1|\le |w|e^{|w|}$ give the second error. The principal scalar
powers remain bounded during this last comparison. This proves the
claim without introducing a factor $\delta^{-2q}$.
\end{proof}

\begin{corollary}[Finite-dimensional scales, zero blocks, and short blocks]
\label{cor:finite-dimensional-details}
For fixed distinct times $0=t_0<t_1<\cdots<t_q$ and fixed vectors
$v_j$, put
$m_j=\lfloor nt_j\rfloor-\lfloor nt_{j-1}\rfloor$,
$z_j=v_j/\sqrt n$, and $\delta_n=\tfrac14n^{-1/32}$.
The spectral product error is bounded by
\[
 C_q\left\{\delta_n^{-4}n^{-1/2}
             +\delta_n^{-6}n^{-1/2}
             +\delta_n^{-2}\sum_j\rho^{m_j}\right\}=o(1).
\]
Repeated times give zero blocks, which may be omitted exactly. A fixed
number of additional deterministic increments of lengths $m_j=o(n)$
may instead be removed from a real characteristic exponent at a cost
$C\sum_j |v_j|\sqrt{m_j/n}$ for the original collision observable.
The estimates are uniform for initial probability densities with uniformly
bounded supremum and variation norms. The short-increment estimate uses
only the supremum bound.
\end{corollary}
\begin{proof}
Distinct fixed times make all $m_j$ positive multiples of $n$ for
large $n$. Substitute the scales in
\eqref{eq:chronological-product-bound}. Smoothing the initial density
costs $O(\delta_n)$; smoothing the covariance costs
$O(\delta_n\ell_{\delta_n})$ by
\eqref{eq:collision-covariance-limit}. The residual-path bound
\eqref{eq:residual-process-maximal} removes smoothing from the original
finite-dimensional process. Integer parts in the covariance time
coefficients cost $O(n^{-1})$. This gives precisely the independent
Gaussian increments in Theorem~\ref{thm:collision-functional}.

When an interval has length zero its operator is the identity, not a
principal spectral projector, so it is omitted before this expansion.
For a deterministic short interval, uniform collision correlation
summability bounds the second moment of its original centered sum by
$Cm_j$, from every deterministic starting index under a bounded
initial density. The difference of the two characteristic exponents,
with the other increments left at their original times, is therefore
bounded in $L^1$ by $C|v_j|\sqrt{m_j/n}$. The elementary exponential
difference bound proves the last assertion. This argument does not
use a complementary spectral power for a short block, and does not
move the other increments in physical time.
\end{proof}

\begin{lemma}[The complex bound before the fourth derivative]
\label{lem:complex-fourth-details}
For one real component $u$ of $h_{R,\delta}$, the stationary pairing
has, on a complex disk $|z|<c\delta^2$, a decomposition
\[
 \int e^{izS_{m,R}u}\dd\nu=e^{m\psi(z)}A(z)+E_m(z)
\]
with $|A(z)|+|\psi(z)|\le C$ and $|E_m(z)|\le C\rho^m$, uniformly
on a slightly smaller such disk. Consequently
$|E_m^{(4)}(0)|\le C\delta^{-8}\rho^m$.
\end{lemma}
\begin{proof}
The two fixed resolvent contours in
Lemma~\ref{lem:smooth-collision-expansion} give uniformly bounded
spectral projections and complementary powers throughout the complex
ball, not just on its real diameter. Both the vector $\nu$ and the
mass functional have bounded norms. Evaluating the complementary
power between them gives the displayed complex bound for $E_m$.
The eigenvalue lies in a fixed disk about one not meeting zero, so its
chosen logarithm and the principal amplitude are uniformly bounded.
Cauchy's estimate on a circle of radius proportional to $\delta^2$
now yields the fourth-derivative estimate. There is no need to bound
the principal power uniformly for complex $z$ and arbitrary $m$.
Its derivative is calculated by the product rule, as in
Lemma~\ref{lem:smooth-fourth-moment}.
\end{proof}

\begin{proposition}[Coarse polygonal interpolation]
\label{prop:coarse-interpolation}
Let $(u_k)$ be bounded real-vector random variables, $|u_k|\le M$.
Let $L\ge1$ be an integer and suppose, under the probability in question,
that every deterministic interval of $m\ge L$ indices has
\[
 \mathbb E\left|\sum_{k=r}^{r+m-1}u_k\right|^4\le C m^2.
\]
At times $jL/n$, interpolate the sums normalized by $\sqrt n$ to obtain
$Z_n$. Let $X_n$ be the polygonal interpolation using every index.
Then, with constants independent of $L,n$,
\begin{equation}\label{eq:coarse-interpolation-bound}
 \mathbb E|Z_n(t)-Z_n(s)|^4\le C'|t-s|^2,
 \qquad \|Z_n-X_n\|_\infty\le 2ML/\sqrt n.
\end{equation}
A final incomplete coarse interval is defined by extending one full
block beyond the observation horizon and restricting the interpolation.
\end{proposition}
\begin{proof}
Write $\Delta=L/n$. For a segment of duration $h\le\Delta$ inside
one block, linear interpolation multiplies the full-block increment by
$h/\Delta$. Its fourth moment is at most
$C\Delta^2(h/\Delta)^4\le Ch^2$. For an interval spanning two adjacent
blocks, apply this estimate separately to the durations $h_1,h_2$ and
use the $L^4$ triangle inequality. Since
$(\sqrt{h_1}+\sqrt{h_2})^4\le4(h_1+h_2)^2$, the same form holds.

For a general interval, split it into its first fractional block, its
complete-block middle interval, and its last fractional block. The
middle interval, when nonempty, contains at least $L$ indices, so the
hypothesis gives the required fourth moment for its total duration.
The $L^4$ triangle inequality and
$(\sqrt{h_1}+\sqrt{h_2}+\sqrt{h_3})^4
 \le9(h_1+h_2+h_3)^2$
prove the first bound with one fixed constant. Empty pieces are simply
omitted. Within one coarse block both interpolations remain within
$ML/\sqrt n$ of the sum at the left endpoint, proving the deterministic
second bound. The endpoint convention uses a full block, so these same
estimates apply on a fixed observation interval without a new constant.
\end{proof}

Applying this proposition with $u_k=h_{R,\delta_n}\circ T_R^k$ and
$L_n=\lceil\delta_n^{-8}\rceil$ gives the interpolation estimates in
Theorem~\ref{thm:collision-functional}. Indeed
\eqref{eq:smooth-fourth-moment} is bounded by $Cm^2$ for $m\ge L_n$.
Stationarity gives the interval estimates under $\nu$, and bounded
initial densities give them by domination. The fourth-moment constants
are uniform. Together with the residual maximal estimate this proves
tightness of the original process, not just of its coarse approximation.

\begin{lemma}[Continuity of the covariance square root]
\label{lem:covariance-square-root}
On every bounded set of positive semidefinite matrices, the map
$C\mapsto C^{1/2}$ is continuous in operator norm. Consequently
$R\mapsto\operatorname{Law}(\mathcal W_{\Gamma_R})$ is continuous on
continuous path space, including at singular matrices.
\end{lemma}
\begin{proof}
Uniformly approximate the scalar square-root function on a bounded
interval containing all the spectra by polynomials. The spectral theorem
bounds the matrix approximation by the same uniform scalar error, and
each polynomial in the matrix is norm-continuous. This proves the first
claim. On one probability space take a standard four-dimensional
Brownian motion $W$ and set $\mathcal W_C=C^{1/2}W$. On any compact time
interval its sample paths are bounded almost surely. Norm continuity
of the square root therefore gives almost-sure uniform convergence of
the coupled paths and hence convergence of their laws.
\end{proof}
